\begin{frame}[fragile]{Adapter Implementation: AutoGen Example}
  \vspace{0.1cm}
  \fontsize{5.5}{6.5}\selectfont
  \begin{block}{}
    \fontsize{5.5}{6.5}\selectfont
    A framework adapter wraps a native agent to conform to the three-method
    OpenCollage contract. Below is a production-grade AutoGen adapter.
  \end{block}

  \vspace{-0.1cm}

  \begin{columns}[T]
    \begin{column}{0.52\textwidth}
      \begin{lstlisting}[style=pythonstyle,basicstyle=\ttfamily\fontsize{3.3}{4}\selectfont,title={\tiny\bfseries opencollage/adapters/autogen\_adapter.py}]
from opencollage.sdk import (
    AdapterBase, AgentConfig, HandoffPayload,
    ExecutionResult, Checkpoint, SchemaRegistry
)
from opencollage.telemetry import OTelEmitter
from autogen import AssistantAgent, UserProxyAgent
from opentelemetry import trace

tracer = trace.get_tracer("opencollage.autogen")

class AutoGenAdapter(AdapterBase):

    def initialize(
        self, config: AgentConfig
    ) -> None:
        self._schema_reg = SchemaRegistry(
            config.schema_registry_url
        )
        self._agent = AssistantAgent(
            name=config.agent_id,
            model=config.model,
            system_message=config.system_prompt,
        )
        self._proxy = UserProxyAgent(
            name=f"{config.agent_id}_exec",
            code_execution_config={
                "executor": config.code_execution
                    .executor,
                "sandbox": config.code_execution
                    .sandbox,
            },
            human_input_mode="NEVER",
        )
        self._output_schema = (
            self._schema_reg.resolve(
                config.output_schema_ref
            )
        )
      \end{lstlisting}
    \end{column}
    \begin{column}{0.48\textwidth}
      \begin{lstlisting}[style=pythonstyle,basicstyle=\ttfamily\fontsize{3.3}{4}\selectfont,title={\tiny\bfseries execute() and emit\_telemetry()}]
    async def execute(self,
        input: HandoffPayload,
        checkpoint: Checkpoint,
    ) -> ExecutionResult:
        with tracer.start_as_current_span(
            f"agent.{self._agent.name}"
        ) as span:
            span.set_attribute("o9e.pipeline_run_id",
                checkpoint.pipeline_run_id)
            chat = await self._proxy.a_initiate_chat(
                self._agent,
                message=input.to_prompt(),
                max_turns=5,
            )
            output = self._extract_output(
                chat.chat_history)
            self._output_schema.validate(output)
            self.emit_telemetry(span, chat)
            return ExecutionResult(
                payload=HandoffPayload(output),
                tokens_used=chat.cost,
                tool_calls=chat.tool_calls,
            )

    def emit_telemetry(self, span, chat):
        span.set_attribute("llm.token.prompt",
            chat.cost["prompt_tokens"])
        span.set_attribute("llm.token.completion",
            chat.cost["completion_tokens"])
        for tc in chat.tool_calls:
            span.add_event("tool_call",
                {"tool": tc.name,
                 "duration_ms": tc.duration})
      \end{lstlisting}
    \end{column}
  \end{columns}
\end{frame}

\begin{frame}[fragile]{Adapter Contract: The Three-Method Interface}
  \scriptsize
  \begin{center}
    \begin{tikzpicture}[scale=0.6, transform shape,
        method/.style={rectangle, draw, thick, rounded corners, fill=orchgreen!15,
                       text width=3.2cm, text centered, minimum height=1.8cm, font=\tiny\bfseries},
        runtime/.style={rectangle, draw, very thick, rounded corners, fill=orchblue!10,
                        text width=11cm, text centered, minimum height=0.8cm, font=\scriptsize\bfseries}]

      \node[runtime] (rt) at (5,3) {\faCogs\ \ OpenCollage Runtime};

      \node[method] (m1) at (0,0) {
        \texttt{initialize(config)}\\[4pt]
        \normalfont\tiny
        Parse AgentConfig YAML\\
        Bootstrap framework agent\\
        Connect to schema registry\\
        Return: \texttt{void}
      };
      \node[method] (m2) at (5,0) {
        \texttt{execute(input, ckpt)}\\[4pt]
        \normalfont\tiny
        Receive typed HandoffPayload\\
        Run native agent logic\\
        Validate output vs schema\\
        Return: \texttt{ExecutionResult}
      };
      \node[method] (m3) at (10,0) {
        \texttt{emit\_telemetry(span)}\\[4pt]
        \normalfont\tiny
        Report tokens, latency, cost\\
        Record tool call events\\
        Propagate trace context\\
        Return: \texttt{void}
      };

      \draw[->, thick, orchblue] (rt.south -| m1) -- (m1);
      \draw[->, thick, orchblue] (rt.south) -- (m2);
      \draw[->, thick, orchblue] (rt.south -| m3) -- (m3);

      \node[font=\tiny\bfseries, orchgray] at (5,-1.6) {
        Every adapter---AutoGen, CrewAI, LangGraph, custom---implements exactly these three methods.
      };
    \end{tikzpicture}
  \end{center}

  \vspace{0.1cm}

  \begin{columns}[T]
    \begin{column}{0.48\textwidth}
      \begin{exampleblock}{\scriptsize Lifecycle Guarantees}
        \fontsize{5.5}{6.5}\selectfont
        \begin{enumerate}
          \setlength{\itemsep}{0pt}
          \item \texttt{initialize} called exactly once at pipeline startup
          \item \texttt{execute} may be called multiple times (retries)
          \item \texttt{emit\_telemetry} called within \texttt{execute}'s span context
          \item Adapter is stateless between \texttt{execute} calls---state lives in checkpoints
        \end{enumerate}
      \end{exampleblock}
    \end{column}
    \begin{column}{0.48\textwidth}
      \begin{alertblock}{\scriptsize Conformance Testing}
        \fontsize{5.5}{6.5}\selectfont
        \begin{itemize}
          \setlength{\itemsep}{0pt}
          \item \texttt{o9e conformance test autogen-adapter}
          \item Validates: schema compliance, telemetry emission, error propagation, idempotent retries
          \item Adapters must pass conformance to be listed in the OpenCollage registry
        \end{itemize}
      \end{alertblock}
    \end{column}
  \end{columns}
\end{frame}
